В современных условиях ускоренного роста цифровых сервисов обеспечение
их надёжности и производительности становится одной из ключевых задач разработки
программного обеспечения.
Нагрузочное тестирование позволяет выявить узкие места системы до её выхода
в промышленную эксплуатацию, оценить поведение сервиса при пиковых нагрузках
и проверить соответствие нефункциональным требованиям.
Среди широко применяемых инструментов --- Apache JMeter, Gatling, k6, Vegeta, hey.
Однако большинство из них требуют тяжёлых зависимостей,
ориентированы исключительно на HTTP или лишены встроенного веб-интерфейса
для визуального мониторинга теста в реальном времени.

\textbf{Актуальность} работы обусловлена потребностью в лёгком инструменте
нагрузочного тестирования, распространяемом в виде единого скомпилированного
бинарного файла, способном работать с HTTP и gRPC сервисами,
интегрироваться в конвейеры CI/CD через REST~API
и предоставлять наглядный веб-интерфейс для контроля хода тестирования.

\textbf{Объект исследования} --- процесс нагрузочного тестирования веб-сервисов
и gRPC-интерфейсов.

\textbf{Предмет исследования} --- программные методы разработки CLI
и веб-интерфейса инструмента нагрузочного тестирования с поддержкой HTTP и gRPC.

\textbf{Теоретическая значимость} работы заключается в систематизации подходов
к проектированию высокопроизводительных CLI-инструментов на языке Go,
анализе механизмов управления параллельными запросами
и ограничение скорости, а также в обосновании архитектурных решений
для потоковой передачи метрик через WebSocket.

\textbf{Практическая значимость} работы состоит в разработке готового к применению
инструмента Tsunami, который позволяет:
нагружать HTTP и gRPC сервисы с настраиваемой частотой запросов;
отслеживать метрики
в режиме реального времени через браузерный интерфейс;
экспортировать результаты в формате JSON для последующего анализа;
интегрировать тестирование в автоматизированные конвейеры через REST~API.
Разработанный инструмент распространяется в виде открытого программного
обеспечения для шести платформ через GitHub Releases и Homebrew,
что обеспечивает его доступность для широкого круга специалистов:
инженеров по качеству, DevOps-специалистов и преподавателей
дисциплин по тестированию программного обеспечения.
Кроме того, проект может служить учебным примером применения
Go-технологий (горутины, WebSocket, gRPC) при подготовке
специалистов в области разработки программного обеспечения.

\textbf{Цель} выпускной квалификационной работы --- разработать CLI-инструмент
и веб-интерфейс для системы нагрузочного тестирования HTTP и gRPC сервисов,
обеспечивающих гибкое управление нагрузкой и визуализацию метрик в реальном времени.

Для достижения поставленной цели необходимо решить следующие \textbf{задачи}:
\begin{enumerate}
  \item Провести сравнительный анализ существующих инструментов нагрузочного тестирования.
  \item Сформулировать функциональные и нефункциональные требования к разрабатываемому инструменту.
  \item Спроектировать архитектуру системы: движок нагрузочного тестирования, CLI-слой и серверную часть.
  \item Реализовать модуль нагрузочного тестирования с поддержкой протоколов HTTP и gRPC.
  \item Реализовать CLI-интерфейс для настройки, запуска и остановки нагрузочных тестов, а также вывода результатов тестирования в терминал.
  \item Разработать веб-интерфейс для запуска тестов и визуализации метрик в реальном времени.
  \item Провести тестирование разработанного инструмента и оценить его производительность.
\end{enumerate}

\textbf{Результатом} бакалаврской выпускной квалификационной работы является
инструмент нагрузочного тестирования Tsunami~--- программное средство
с интерфейсом командной строки и встроенным веб-приложением,
реализующее нагрузочное тестирование HTTP- и gRPC-сервисов
с отображением метрик в режиме реального времени.

\textbf{Структура бакалаврской выпускной квалификационной работы.}
Выпускная квалификационная работа состоит из введения, двух глав, заключения,
списка литературы и приложений.
В первой главе проводится анализ предметной области:
рассматриваются существующие инструменты нагрузочного тестирования,
обосновывается выбор технологий и проектируется архитектура системы.
Во второй главе описывается практическая реализация:
движок нагрузочного тестирования, CLI-интерфейс, серверная часть
с поддержкой WebSocket и веб-приложение на React.
